The following two lemmas provide key ingredients for the proof of Theorem~\ref{main theorem prep}. For $k\in [n]$, we let $y_k$ denote $x_k+x_{-k}$.

\begin{lemma}\label{lm: symmetric stress expression}
Let $w \in \R[X]_i$ be a squarefree symmetric polynomial such that $\partial_{x_v} w$ is symmetric for all vertices $v$. Then $w$ is a squarefree polynomial in $y_1,\ldots,y_n$, that is, $w$ can be written as $$w=\sum_{\substack{\tau\subseteq [n]\\|\tau|=i}} c_\tau \prod_{k\in \tau} (x_k+x_{-k}) \quad \mbox{for some $c_\tau\in \R$}.$$ 
\end{lemma}
\begin{proof}
 It is easy to prove by induction on $n$ that a squarefree polynomial $Q\in\R[X]$ is a polynomial in $y_1,\ldots,y_n$ if and only if $\partial_{x_{k}}Q=\partial_{x_{-k}}Q$ for all $k\in [n]$. Thus 
to prove the lemma, it is enough to check that our given $w$ satisfies $\partial_{x_{k}}w=\partial_{x_{-k}}w$ for all $k\in [n]$. Indeed, by symmetry of $w$ and $\partial_{x_k} w$, and by the definition of $\alpha$,
\[ 
\partial_{x_{k}}w=\alpha(\partial_{x_{k}}w)=\partial_{x_{-k}} (\alpha w)=\partial_{x_{-k}} w.
\]
The result follows.
\end{proof}


\begin{lemma}\label{lm: i-stress to i+1-stress}
Let $i\geq 1$ and let $w\in \R[X]_{i+1}$ be a squarefree polynomial such that for all vertices $v$, $\partial_{x_v} w$ is a polynomial in $y_1,\ldots,y_n$. Then $w$ is a squarefree polynomial in $y_1,\ldots,y_n$. In particular, $w$ is symmetric and can be expressed as $$w=\sum_{\substack{\sigma\subseteq [n]\\|\sigma|=i+1}} c_\sigma \prod_{k\in \sigma} (x_k+x_{-k}) \quad \mbox{for some }c_\sigma\in\R.$$
\end{lemma}
\begin{proof}
By Lemma~\ref{lm: symmetric stress expression}, the statement will follow if we show that $w$ is symmetric. To check this, write $w$ as $w=\sum c_{k_1,k_2,\ldots,k_{i+1}} x_{k_1}x_{k_2}\cdots x_{k_{i+1}}$ for some $c_{k_1,k_2,\ldots,k_{i+1}}\in \R$. The assumption that partial derivatives of $w$ are polynomials in $y_1,\ldots,y_n$ implies that $\partial_{x_{k_2}}\cdots \partial_{x_{k_{i+1}}} w$ is symmetric. Hence $c_{k_1,k_2,\ldots,k_{i+1}}= c_{-k_1,k_2,\ldots,k_{i+1}}$ (as they are coefficients of $x_{k_1}$ and $x_{-k_1}$ in $\partial_{x_{k_2}}\cdots \partial_{x_{k_{i+1}}} w$). Repeated applications of this argument imply that $c_{k_1,k_2,\ldots,k_{i+1}}= c_{-k_1,-k_2,\ldots,-k_{i+1}}$. Thus, $w$ is symmetric.
\end{proof}
